\appendix
\chapter{Annexes}
\markboth{Annexes}{}

\section{Schéma unifié des données}
Le tableau~\ref{tab:schema} décrit les principales colonnes du jeu de données
unifié.

\begin{table}[H]
\centering
\caption{Principales colonnes du schéma unifié.}
\label{tab:schema}
\begin{tabular}{ll}
\toprule
\textbf{Colonne} & \textbf{Description} \\
\midrule
Societe & Société émettrice \\
Statement & État financier (bilan, résultat, flux, SIG) \\
Indicateur & Libellé lisible du poste \\
Cle\_indicateur & Clé normalisée (jointure des séries) \\
Date / Annee / Periode & Repères temporels (T1, S1, 9M, FY) \\
Valeur & Montant \\
Devise / Echelle & TND/USD ; unité/million \\
Type\_donnee & Réel ou estimé \\
\bottomrule
\end{tabular}
\end{table}

\section{Extrait de code : méthode d'augmentation}
L'extrait suivant illustre la logique de génération des grandeurs de flux
(répartition au prorata du temps).

\begin{lstlisting}[language=Python,caption={Augmentation des grandeurs de flux (extrait).}]
# Flux cumules : T1 = moitie du semestre ; 9M = S1 + moitie du 2e semestre
if jun in val and val[jun] is not None:
    nouvelles.append(_ligne(modele, mar, "T1", 0.5 * val[jun]))
if jun in val and dec in val:
    nouvelles.append(_ligne(modele, sep, "9M",
                            val[jun] + 0.5 * (val[dec] - val[jun])))
\end{lstlisting}

\section{Extrait de code : exemple de ratio et de notation}
\begin{lstlisting}[language=Python,caption={Définition d'un ratio et d'une grille de notation (extrait).}]
# Ratio de liquidite generale = Actifs courants / Passifs courants
"LIQ_GEN": lambda r: _d(r.actif_courant, r.dettes_courantes)

# Grille de notation 0-5 (benchmark 1,5 a 2,0)
"LIQ_GEN": [(None,1.0,0),(1.0,1.2,1),(1.2,1.5,2),(1.5,2.0,4),(2.0,None,5)]
\end{lstlisting}

\section{Requêtes SQL d'exemple (entrepôt)}
\begin{lstlisting}[language=SQL,caption={Requêtes analytiques sur l'entrepôt.}]
-- Score SGPI par societe (exercice 2024)
SELECT societe, sgpi_sur_100 FROM v_sgpi
WHERE annee = 2024 AND periode = 'FY' ORDER BY sgpi_sur_100 DESC;

-- Benchmark de la marge nette
SELECT societe, valeur FROM v_ratios
WHERE code_ratio='MARGE_NETTE' AND annee=2024 AND periode='FY'
ORDER BY valeur DESC;
\end{lstlisting}
